% ─────────────────────────────────────────────────────────────────────────────
\section{Supporting Standards and Software References}
\label{sec:appendix-supporting-reference-framework}
% ─────────────────────────────────────────────────────────────────────────────

The experimental and reproducibility framework is supported by additional
references concerning general-purpose machine-learning software, vendor-specific
media-classification capabilities, forensic methodology, and tool validation.
The standard reference for scikit-learn supports the discussion of conventional
supervised-learning workflows and evaluation utilities presented in
Chapter~\ref{chap:background} \cite{pedregosa2011scikit}. The description of
automated media-classification capabilities associated with the Cellebrite
software perimeter is further grounded in the corresponding vendor documentation
\cite{cellebrite_media_classification}.

From a methodological perspective, ISO/IEC 27041 addresses the suitability and
adequacy of investigative methods, ISO/IEC 27042 concerns the analysis and
interpretation of digital evidence, and ISO/IEC 27043 provides principles and
processes for digital incident investigation
\cite{iso27041,iso27042,iso27043}. These standards complement the evidence-handling
principles discussed throughout the thesis. The validation controls applied to
the experimental artifact are also consistent with the broader requirement that
forensic methods and tools undergo documented validation before their results are
relied upon in operational contexts \cite{swgde2014validation}.
